\section{Attack Narrative}

\subsubsection*{Introduction}
\textcolor{red}{Intro here if needed!!}

\subsection{Methodology}

\team utilized the Penetration Testing Execution Standard (PTES) to conduct a comprehensive evaluation of \clientP network security. This methodology offers a structured approach that assesses the organization's security posture from both technical and business perspectives, ensuring that the security measures align with \clientP objectives and safeguard customer well-being.

The PTES framework encompasses seven critical phases:

\begin{enumerate}
    \item \textbf{Pre-Engagement Interactions}
    \item \textbf{Intelligence Gathering}
    \item \textbf{Threat Modeling}
    \item \textbf{Vulnerability Analysis}
    \item \textbf{Exploitation}
    \item \textbf{Post-Exploitation}
    \item \textbf{Reporting}
\end{enumerate}

Each phase builds upon the previous one to systematically identify and address potential security vulnerabilities. By adhering to this standard, \team ensured a thorough and methodical assessment of the \client network.

For a detailed explanation of each phase and how it was applied during this engagement, please refer to \hyperref[sec:appendixE]{Appendix E}.

\newpage

\input{content/engagementOverview/phishingAssessment}

\